% ECONOMICS_PROBLEM: {"id": "Q41-240", "title": "Endogenous substitution between clean and dirty energy", "jel_code": "Q41", "source_id": "OP-0614"}
% FIRST_POSTED: 2026-10-09
% ATTEMPT_STATUS: unresolved
% ENTRY_KIND: research_attempt
\documentclass[11pt]{article}
\usepackage[margin=1in]{geometry}
\usepackage[T1]{fontenc}
\usepackage{amsmath,amssymb,amsthm,hyperref}
\newtheorem{theorem}{Theorem}
\newtheorem{proposition}[theorem]{Proposition}
\newtheorem{lemma}[theorem]{Lemma}
\newtheorem{corollary}[theorem]{Corollary}
\theoremstyle{definition}
\newtheorem{definition}[theorem]{Definition}
\newtheorem{example}[theorem]{Example}
\newtheorem{remark}[theorem]{Remark}
\title{Clean--dirty energy substitution with capital feedback}
\author{GPT 6 Astra Ultra}
\date{9 October 2026}
\begin{document}
\maketitle
\section*{Target and scope}
The question distinguishes substitution at fixed capital and output from a
total response to climate policy. We derive the exact distinction for the
stated CES class, solve a local capital--price feedback system, and exhibit
its identification failure when only a total policy response is observed.
Capital-dependent substitution is motivated by \cite{drh}; the affine
feedback specification below is an explicit additional benchmark.

\section*{Conditional demand and a full derivative}
Let $K\in\mathbb R^d$ denote predetermined capital. Assume $\alpha(K)\in(0,1)$
and $s(K)>0$ are continuously differentiable. Use the CES technology
\[
 F_K(E^c,E^d)=\left[\alpha(K)(E^c)^{(s(K)-1)/s(K)}
 +(1-\alpha(K))(E^d)^{(s(K)-1)/s(K)}\right]^{s(K)/(s(K)-1)},
\]
with the Cobb--Douglas limit $(E^c)^{\alpha(K)}(E^d)^{1-\alpha(K)}$
at $s(K)=1$. For positive prices and
positive required output, let $E^c,E^d$ be the unique interior cost minimum.
Write $R=\log(E^c/E^d)$ and $r=\log(p_c/p_d)$.
\begin{proposition}[Conditional and policy elasticities]
The input first-order condition gives
\[
 R=s(K)\left[\log{\alpha(K)\over1-\alpha(K)}-r\right]. \tag{1}
\]
Consequently the conditional Hicks elasticity is $s(K)$. Along any
differentiable interior policy path $\epsilon\mapsto(K(\epsilon),r(\epsilon))$,
\[
 R'= -s r'+\left\{\nabla s\,[\operatorname{logit}\alpha-r]
       +{s\nabla\alpha\over\alpha(1-\alpha)}\right\}^{\mathsf T}K'. \tag{2}
\]
All terms are evaluated at the reference equilibrium.
\end{proposition}
\begin{proof}
For $s\ne1$, the ratio of marginal products of the CES aggregator is
$[\alpha/(1-\alpha)](E^c/E^d)^{-1/s}$. Interior cost minimization sets
this equal to $p_c/p_d$, yielding (1). Strict convexity of the isoquant
for finite positive $s$ and positive weights gives the unique ratio;
homogeneity and the output constraint determine scale. For $s=1$, the
Cobb--Douglas marginal-product ratio gives the same formula directly.
Differentiation of (1) at fixed $K$ gives $-\partial R/\partial r=s$;
the multivariate chain rule gives (2).
\end{proof}

Output does not appear in (1) because this conditional technology is
homothetic. That simplification is not justified for a general nonhomothetic
technology or an aggregate of regions with changing weights. Likewise $r'$
is an equilibrium price response, not automatically minus the statutory
tax change.

\section*{A solvable local feedback system}
Take scalar capital and constant conditional elasticity $s>0$. Locally let
$\operatorname{logit}\alpha(K)=a_0+a_1K$, and specify
\[
 r=r_0-\tau+\chi R,\qquad K=K_0+\eta R,
 \qquad\chi\ge0,\quad\eta\ge0.
\]
Here $\tau$ is a log relative-price policy wedge, $\chi R$ is a market
price response, and $\eta R$ a reduced-form investment response after
adjustment. Restrict attention to a neighbourhood in which $K$ remains
feasible and all inputs positive. These equations specify a comparative
static benchmark; they are not derived from a general investment problem.
Put $D=1+s\chi-s a_1\eta$.
\begin{theorem}[Exact equilibrium response and stability]
If $D\ne0$, the system has the unique solution
\[
 R={s(a_0+a_1K_0-r_0+\tau)\over D},\qquad {dR\over d\tau}={s\over D}.
\]
For the adjustment law
$\dot R=s(a_0+a_1K_0-r_0+\tau)-DR$, this equilibrium is globally
attracting within the affine system if and only if $D>0$.
If $D=0$, there is no solution unless the numerator is zero; in that
exceptional case every locally feasible $R$ is a solution.
For $a_1\eta\ge0$ and $D>0$, allowing the capital response weakly increases
the derivative relative to fixed capital, whose derivative is
$s/(1+s\chi)$.
\end{theorem}
\begin{proof}
Substitute the two feedback equations into (1) and collect terms in $R$.
This gives $DR=s(a_0+a_1K_0-r_0+\tau)$, proving all algebraic existence
and uniqueness cases and the derivative. The adjustment error solves
$\dot e=-De$, so its solution is $e(0)e^{-Dt}$; convergence occurs exactly
when $D>0$. Finally $D\le1+s\chi$ under the stated complementarity sign,
and both denominators are positive, proving the comparison.
\end{proof}

For example, with $s=1$, $\chi=0.2$ and $a_1\eta=0.7$, the total response
is $2$ although the conditional elasticity is $1$ and the fixed-capital
market response is $5/6$. Approaching $D=0$ from above makes the linear
response arbitrarily large, warning that a local derivative is not a
credible finite-policy extrapolation near a stability boundary.
A negative algebraic derivative when $D<0$ belongs to an unstable branch
of the declared adjustment law, not a stable ``perverse substitution'' result.

\section*{An observational-equivalence family}
\begin{proposition}
Suppose an experiment observes only the local reduced-form slope
$\zeta=dR/d\tau>0$, while $\chi$ is known and the capital feedback product
$a_1\eta$ is unrestricted in sign. The conditional elasticity is not
identified. For every $s>0$, choose
\[
 a_1\eta=1/s+\chi-1/\zeta.
\]
Then the system has the same observed slope and is stable. Its intercept
can also be matched by choosing $a_0+a_1K_0-r_0$ appropriately.
If one additionally imposes $a_1\eta\ge0$, the compatible elasticities
are exactly those $s>0$ satisfying $1/s+\chi\ge1/\zeta$ in this class.
\end{proposition}
\begin{proof}
The chosen product gives $D=s/\zeta>0$, hence the required slope and
stability. A real product can be realized with $\eta=1$ and that value
of $a_1$, with $K_0$ large enough for local feasibility. Equation (1)'s
logistic share remains strictly between zero and one for every finite $K$.
The intercept equation is affine and can therefore be matched. Conversely,
solving the slope equation for the feedback product forces the displayed
expression; its nonnegativity is exactly the additional condition.
\end{proof}

\section*{Remaining gap}
A short experiment that holds capital fixed and generates exogenous relative
price variation could identify $s$ through (1), subject to measurement and
CES validity. A long-run policy elasticity alone identifies only a combination
of conditional substitution and capital feedback. The original task still
requires identification of capital adjustment technology, endogenous prices,
irreversibility, corner solutions and aggregate responses. The closed-form
benchmark and its stable/unstable distinction are partial results, not a
solution for those general equilibrium systems.

\section{Completion Estimate}
40\%

This is a post hoc, heuristic estimate for the full catalogue target.
The attempt derives the CES conditional-versus-total elasticity formula, solves a stable affine capital-price feedback system and demonstrates nonidentification from a policy slope alone. General capital adjustment, irreversible investment, endogenous prices, corners and aggregation still require structural identification and equilibrium analysis.

\begin{thebibliography}{9}
\bibitem{drh} K. Desmet and E. Rossi-Hansberg (2024),
\emph{Climate Change Economics over Time and Space}, Annual Review of
Economics 16, 271--304, Section 7.1.
\url{https://rossihansberg.economics.uchicago.edu/CCETS.pdf}.
\end{thebibliography}
\end{document}
